\input{preamble}

\begin{document}

\title{Groups}
\maketitle

\phantomsection
\label{section-phantom}

\tableofcontents

\section{Group extensions and semidirect
products}
\label{section-group-extensions}

\begin{definition}
\label{definition-group-extension}
A {\it group extension} of a group $Q$ by a
group $N$ is a short exact sequence
$$
1\longrightarrow N
\overset{i}{\longrightarrow}G
\overset{\pi}{\longrightarrow}Q
\longrightarrow1.
$$
It is {\it split} if there is a group
homomorphism
$$
s:Q\longrightarrow G
$$
such that
$$
\pi\circ s=\operatorname{id}_Q.
$$
Such a map $s$ is called a {\it splitting}
or a {\it section}.
\end{definition}

\noindent
Exactness says in particular that
$i(N)=\operatorname{Ker}\pi$, so $i(N)$ is
normal in $G$. We usually identify $N$ with
its image in $G$.

\begin{proposition}
\label{proposition-split-extension-semidirect}
If
$$
1\longrightarrow N
\longrightarrow G
\longrightarrow Q
\longrightarrow1
$$
splits, then
$$
G\cong N\rtimes_\varphi Q,
$$
where
$$
\varphi:Q\longrightarrow
\operatorname{Aut}(N)
$$
is defined by
$$
\varphi(q)(n)=s(q)ns(q)^{-1}.
$$
\end{proposition}

\begin{proof}
Let $g\in G$ and put $q=\pi(g)$. Then
$$
\pi(gs(q)^{-1})=1,
$$
so there is $n\in N$ such that
$$
g=ns(q).
$$
This expression is unique: if
$$
ns(q)=n's(q'),
$$
applying $\pi$ gives $q=q'$, and then
$n=n'$.

Thus the map
$$
N\times Q\longrightarrow G,
\qquad
(n,q)\longmapsto ns(q)
$$
is a bijection. Multiplication in these
coordinates is
$$
(n,q)(n',q')
=
(n\varphi(q)(n'),qq'),
$$
which is the multiplication in the
semidirect product.
\end{proof}

\begin{remark}
\label{remark-split-not-direct}
A split extension need not be a direct
product. It is a direct product with respect
to the chosen splitting precisely when
$$
\varphi(q)=\operatorname{id}_N
$$
for every $q\in Q$, equivalently when
$s(Q)$ centralizes $N$.
\end{remark}

\begin{example}
\label{example-s3-semidirect}
The sequence
$$
1\longrightarrow C_3
\longrightarrow S_3
\longrightarrow C_2
\longrightarrow1
$$
splits, so
$$
S_3\cong C_3\rtimes C_2.
$$
The nontrivial element of $C_2$ acts on
$C_3$ by inversion. Hence the action is not
trivial, and
$$
S_3\not\cong C_3\times C_2.
$$
\end{example}

\end{document}
